\begin{figure*}[htb]
    \centering
    \includegraphics[width=\textwidth]{Figures/Composite.png}
    \caption{Entangling Gadget. After each operation, we project the oscillator onto a position eigenstate and then compute the expectation values of oscillator-qubit operators $\langle \Pi_x \otimes \sigma_j \rangle / \langle \Pi_x\rangle$ for $j\in\{X,Y,Z\}$. The vectors in the $z$-$y$ plane of the Bloch sphere, corresponding to these expectation values are shown in green, for each position eigenvalue. (a) Initial oscillator-ancilla state with $\Delta=1$. The oscillator is either a coherent state located at $x=\pm\alpha$. (b) Rotation error in ancilla due to quantum uncertainty upon conditional momentum boost by $\frac{\pi}{4|\alpha|}$. (c) Oscillator-ancilla entanglement upon conditional displacement by $\frac{\pi}{4|\alpha|}$. Note that this operation applies a preemptive correction to errors in (b). (d) Oscillator-ancilla state followed by operations (1) and (2). The measurement of the ancilla state at this point correctly identifies the sign of the position eigenvalue at the peak of  $\psi(x)$. This entanglement is subjected to (1) applied to the ancilla in $\ket{g}$, that is, the target unitary (2).}
    \label{fig:Composite}
\end{figure*}

\begin{figure*}[htb]
     \centering
     \includegraphics[width=\textwidth]{Figures/Framework.png}
     %\caption{Entangling and disentangling sequence between continuous variable states in the oscillator and ancilla qubits is shown above. The correction cancels the error in ancilla rotation due to the quantum uncertainty in position eigenvalue $x$. The disentangling sequence is not only adjoint of the entangling operation but also has other options due to the periodicity in ancilla rotation. Condition on $\theta'$ and $\lambda'$ is given by Eq.~\ref{eq:unent}. \added{Add an example like GKP as well.}}
     \label{fig:Framework}
 \end{figure*}